\section*{Abstract}
Multimodal emotion recognition research has predominantly framed the problem as one of classification accuracy, seeking to identify a single ground-truth emotion label from complementary behavioural signals. This thesis reframes the problem twice. First, when an individual's external behavioural signals (audio and video) and their internal neurophysiological signal (EEG) disagree, that disagreement is itself a clinically meaningful quantity, associated in the literature with emotion suppression, masked affect, alexithymia, and affect dysregulation disorders. We term this phenomenon cross-modal affective coherence and develop a computational framework to detect and characterise it. Second, we show that the evaluation protocol under which such systems are conventionally assessed on this benchmark substantially overstates what they achieve on people they have not seen.

The framework is evaluated on the EAV dataset \cite{Lee2024EAV}, a cue-based laboratory corpus of 42 subjects with synchronised 30-channel EEG, audio, and video recordings across five emotion categories. Per-modality encoders---the Audio Spectrogram Transformer (AST), a facial-emotion Vision Transformer (ViT), and EEGNet---are fine-tuned on EAV, and a trimodal cross-modal attention fusion module with 2.28 million parameters is trained on cached per-modality features. Results are reported under two protocols: the conventional within-subject protocol, and a five-fold subject-wise cross-validation in which every participant is held out entirely and tested as an unseen individual.

Under the within-subject protocol, per-modality mean accuracies are audio 57.1\%, vision 74.6\%, and EEG 43.9\%, all substantially exceeding the EAV paper's published baselines of 36.7\%, 52.8\%, and 36.7\%. Multimodal fusion yields 80.2\% (cross-attention), 81.7\% (concat-MLP) and 84.7\% (with softhard modality dropout), against 77.5\% for naive late fusion.

Under subject-independent evaluation, these figures change substantially and informatively. The best system reaches 62.2\%, and the ordering of methods reverses: naive late fusion significantly outperforms every learned fusion variant ($p < 0.003$), and the three learned variants become statistically indistinguishable from one another. The loss is not evenly distributed across modalities. Audio transfers at parity ($+0.4$ pp), EEG loses 5.7 points, and vision collapses by 33.4 points---from the strongest modality to the second weakest. Training the vision encoder longer makes its cross-subject accuracy worse. We interpret this as evidence that EAV's high within-subject vision accuracy is substantially identity-driven rather than expression-driven, with implications for every within-subject result reported on this benchmark.

The coherence contribution is a $5\times5$ suppression matrix computed over 5,040 test trials. Trials qualifying as suppression events---audio-vision consensus with an EEG disagreement at confidence $\geq 0.5$---total 401 (8.0\% base rate). The dominant pattern is a bidirectional Anger--Happiness confusion (117 events), consistent with shared high-arousal physiology in EEG. Secondary directional patterns---Sadness$\rightarrow$Neutral (43) and Calmness$\rightarrow$Neutral (40)---reflect low-arousal emotion suppression. All patterns are cross-population rather than artefacts of individual outliers. This work constitutes the first characterisation of cross-modal affective coherence on the EAV benchmark, and the first subject-independent evaluation of trimodal fusion on it.

\vspace{1cm}
\textbf{Keywords: Multimodal Emotion Recognition, EEG, Affective Coherence, Emotion Suppression, Cross-Modal Attention, Subject-Independent Evaluation, Audio Spectrogram Transformer, Vision Transformer, EEGNet, EAV Dataset, Affective Computing.}
\pagebreak
